\documentstyle[%


righttag%


,11pt%


,amssymb%


,ctagsplt%


,russian%               remove in Eng.


,izvru%                 Set izven% in Eng.


% ,izvru-ks             for brief communications


]{amsart}





%       Math definitions





\newcommand{\trid}{\dots\dots\dots}


\newcommand{\rang}{\operatorname{rang}}


\newcommand{\tts}[1]{}





\theoremstyle{plain}


\theorembodyfont{\it}


\newtheorem{thmnonum}{\hskip\parindent Теорема}


\renewcommand{\thethmnonum}{}





%\hoffset=-15mm                  For Eng.





\setcounter{page}{3}


\god{1997}


\nomer{9~(424)}%                Remove in Eng.


%\nomer{Vol.\,41, No.\,9}        Set in Eng.


% \str{75-81}                    Set in Eng.


\udk{517.73}


\author{\it Л.С.\,Велимирович}


\title{Бесконечно малые изгибания торообразной\\


поверхности вращения с многоугольным меридианом}





\date{}


% \pagestyle{myheadings}         Set in Eng.





\pagestyle{plain} %              Remove in Eng.


\begin{document}





% \thispagestyle{plain}          Set in Eng.





\maketitle





% Body text





{\bf Введение.} В этой работе рассмотрены бесконечно малые (б.\,м.)


изгибания торообразной поверхности вращения с многоугольным меридианом и


получено необходимое и достаточное условие, чтобы такая поверхность была


нежесткой. Также дана процедура определения поля б.\,м. изгибаний для


такой поверхности. Результат К.М.\,Белова \cite{Bel} является частным


случаем теоремы, доказанной в этой работе.





Пусть простой многоугольник $p_n$ с вершинами $A_i(u_i,\rho_i)$


($i=1,\dotsc,n$) в


декартовой системе


координат $u0\rho$ вращается вокруг оси $u$, единичный вектор которой


$\Bar e$. Тогда уравнение сторон имеют вид


\begin{equation}


A_mA_{m+1}:\;\rho_{(m)}=\rho_m+


\frac{\rho_{m+1}-\rho_m}{u_{m+1}-u_m}(u-u_m),\quad


\rho_{(m)}'=\frac{\rho_{m+1}-\rho_m}{u_{m+1}-u_m}=k_m \quad


(m=1,\dotsc,n; \ A_{n+1}\equiv A_1),


\end{equation}


где $\rho_{(m)}$ --- значение $\rho$ на $A_mA_{m+1}$.





Чтобы рассмотреть б.\,м. изгибания торообразной поверхности вращения с


простым замкнутым меридианом, мы используем метод Кон-Фоссена


\cite{Kohn}.





Радиус-вектор точки поверхности будет


\begin{equation}


\Bar r(u,v)=u\Bar e+\rho(u)\Bar a(v),


\end{equation}


где $\rho=\rho(u)$ --- уравнение меридиана, $\Bar a(v)$ --- единичный


вектор оси $\rho$, $v$


--- угол между плоскостью начального положения меридиана и


$\Bar a(v)$, $\Bar e$ ---


единичный вектор оси вращения.





Фундаментальное поле б.\,м. изгибаний поверхности (2) будет


\begin{equation}


\Bar z(u,v)=[\varphi_k(u)e^{ikv}+\Tilde\varphi_ke^{-ikv}]\Bar e


+[\psi_k(u)e^{ikv}+\Tilde\psi_k(u)e^{-ikv}]\Bar a(v)+


[\chi_k(u)e^{ikv}+\chi_k(u)e^{-ikv}]\Bar a'(v),


\end{equation}


где, например, $\Tilde\varphi_k(u)$ --- сопряженное значение для


$\varphi_k(u)$,


функции $\varphi_k(u)$, $\psi_k(u)$, $\chi_k(u)$


удовлетворяют уравнениям


\begin{equation}


\begin{split}


&\varphi_k'(u)-\rho'(u)\psi_k'(u)=0,\quad


\psi_k(u)+ik\chi_k(u)=0,\\


&ik\varphi_k(u)+\rho'(u)[ik\psi_k(u)-\chi_k(u)]+


\rho(u)\chi_k'(u)=0.


\end{split}


\end{equation}


Функции $\psi_k(u)$, $\chi_k(u)$ удовлетворяют также и уравнению


\begin{equation}


\rho(u)\lambda''(u)+(k^2-1)\rho''(u)\lambda(u)=0,


\end{equation}


где $\lambda(u)$ --- неизвестная функция. Опуская индекс $k$, обозначим


через $\psi_{(i)}$ значение $\psi$ на сторонах


$A_iA_{i+1}$, $i=1,\dotsc,n$, $A_{n+1}\equiv A_1$.





Из уравнений (1) и (5) вытекает линейность функций $\psi_{(i)}(u)$,


т.\,е.


\begin{equation}


\psi_{(i)}(u)=M_iu+N_i\quad (i=1,\dotsc,n).


\end{equation}


В точках $u=\sigma$ меридиана, где


$\rho'(\sigma-0)\ne\rho'(\sigma+0)$, т.\,е. в вершинах многоугольника,


предполагая непрерывность функций $\psi_{(i)}(u)$, получаем


\begin{equation}


\psi_{(i)}(u_i)=\psi_{(i-1)}(u_i),\quad i=1,\dotsc,n;\quad


\psi_{(1)}(u_1)=\psi_{(n)}(u_1)


\end{equation}


и отсюда, учитывая (6), имеем


\begin{equation}


M_iu_i+N_i=M_{i-1}u_i+N_{i-1},\quad i=1,\dotsc,n;\quad


M_0\equiv M_n\quad N_0=N_n.


\end{equation}


Рассматривая эту систему как систему с неизвестными $N_i$,


$i=1,\dotsc,n$, получаем


\begin{equation}


\begin{split}


N_1-N_n&=-M_1u_1+M_nu_1,\\


N_1-N_2&=-M_1u_2+M_2u_2,\\


&\dots\\


N_{m-1}-N_m&=-M_{m-1}u_m+M_mu_m,\\


&\dots\\


N_{n-1}-N_n&=-M_{n-1}u_n+M_mu_n.


\end{split}


\end{equation}


В вершинах многоугольника имеем уравнение (\cite{Kohn}, c.\,112,


ур.\,(15))


\begin{equation}


\rho(\sigma)[\psi'(\sigma+0)-\psi'(\sigma-0)]+


(k^2-1)\psi(\sigma)[\rho'(\sigma+0)-\rho'(\sigma-0)]=0,


\end{equation}


откуда получается система


$$


\rho_i(M_i-M_{i-1})+(k^2-1)(M_iu_i+N_i)(k_i-k_{i-1})=0


\quad (i=1,\dotsc,n; \ M_0\equiv M_n, \ k_0\equiv k_n),


$$


т.\,е.


\begin{equation}


\begin{split}


[\rho_1+(k^2-1)(k_1-k_n)u_1]M_1-\rho_1M_n+(k^2-1)(k_1-k_n)N_1&=0,\\


-\rho_2M_1+[\rho_2+(k^2-1)(k_2-k_1)u_2]M_2+


(k^2-1)(k_2-k_1)N_2&=0,\\


&\dots\\


-\rho_mM_{m-1}+[\rho_m+(k^2-1)(k_m-k_{m-1})u_m]M_m


+(k^2-1)(k_m-k_{m-1})N_m&=0,\\


&\dots\\


-\rho_nM_{n-1}+[\rho_n(k^2-1)(k_n-k_{n-1})u_n]M_n


+(k^2-1)(k_n-k_{n-1})N_n&=0.


\end{split}


\end{equation}


Уравнения (9) и (11) составляют систему линейных уравнений с неизвестными


$M_i$, $N_i$\linebreak %%%%%%%% Remove in Eng.


($i=1,\dotsc,n$).





Рассмотрим систему (9) с неизвестными $N_i$.


Для расширенной матрицы этой системы


\begin{equation}


P=[p_{ij}]\quad (i=1,\dotsc,n; \ j=1,\dotsc,n+1)


\end{equation}


имеем


\begin{gather*}


p_{11}=1,\quad p_{1n}=-1,\quad p_{1\,n+1}=-M_1u_1+M_nu_1;\quad


p_{1n}=0,\quad r=2,\dotsc,n-1,\\


p_{m\,m-1}=1,\quad p_{mm}=-1,\quad p_{m\,n+1}=(M_m-M_{m-1})u_m,


\quad p_{ms}=0,\\


s\notin\{m-1,m,n+1\},\quad m=2,\dotsc,n.


\end{gather*}


Элементарными преобразованиями: 1${}^\circ$~($-$I)$\to$(II),


2${}^\circ$~($k-1$)$\to$($k$),


матрицу


(12) приводим к виду


\begin{equation}


P\sim


\begin{bmatrix}


N_1 & N_2 & & N_{m-1} & N_m & & N_{n-1} & N_n\\


1& 0 & \dots & 0 & 0 & \dots & 0 & -1 & p_1\\


0 & -1 & \dots & 0 & 0 & \dots & 0 & 1 & p_2\\


\dots & \dots & \dots & \dots & \dots & \dots & \dots & \dots & \dots \\


0 & 0 & \dots & 0 & -1 & \dots & 0 & 1 & p_m\\


\dots & \dots & \dots & \dots & \dots & \dots & \dots & \dots & \dots \\


0 & 0 & \dots & 0 & 0 & \dots & -1 & 1 & p_{n-1}\\


0 & 0 & \dots & 0 & 0 & \dots & 0 & 0 & p_n


\end{bmatrix}


,


\end{equation}


где


\begin{equation}


\begin{split}


p_1&=(M_n-M_1)u_1,\\


p_2&=(M_1-M_n)u_1+(M_2-M_1)u_2,\\


&\dots\\


p_m&=(M_1-M_n)u_1+\sum^m_{l=1}(M_l-M_{l-1})u_l,\quad


m=2,\dotsc,n,\\


&\dots\\


p_{n-1}&=(M_1-M_n)u_1+\sum^{n-1}_{l=2}(M_l-M_{l-1})u_l,\\


p_n&=(M_1-M_n)u_1+\sum^n_{l=2}(M_l-M_{l-1})u_l.


\end{split}


\tag{$13{}'$}


\end{equation}





Система (9) совместна тогда и только тогда, когда


$\rang M=\rang P$, где $M$ --- матрица


системы (9), $P$ --- расширенная матрица этой системы, т.\,е.


тогда и только тогда, когда


$$


(M_1-M_n)u_1+\sum^n_{l=2}(M_l-M_{l-1})u_l=0,


$$


т.\,е.


\begin{equation}


M_m=\frac{1}{u_1-u_n}\sum^{n-1}_{i=1}(u_i-u_{i-1})M_i.


\end{equation}


Для редуцированной системы (согласно (13) и (13$'$)) имеем


\begin{equation}


\begin{split}


N_1-N_n&=(M_n-M_1)u_1,\\


-N_2+N_n&=(M_1-M_n)u_1+(M_2-M_1)u_2,\\


&\dots\\


-N_m+N_n&=(M_1-M_n)u_1+\sum^m_{l=2}(M_l-M_{l-1})u_l,\\


&\dots\\


-N_{n-1}+N_n&=(M_1-M_n)u_1+\sum^{n-1}_{l=2}(M_l-M_{l-2})u_l.


\end{split}


\end{equation}


Вводя обозначения


$$


u_i-u_j=u_{ij},\quad k_i-k_j=k_{ij},


$$


из (15) получаем


\begin{equation}


\begin{split}


N_1&=N_n+\frac{u_1}{u_{1\,n}}u_{n\,2}M_1+\frac{u_1}{u_{1\,n}}


\sum^{n-1}_{i=2}u_{i\,i+1}M_i,\\


&\dots\\


N_m&=N_n+\frac{u_n}{u_{1\,n}}\sum^{m-1}_{i=1}u_{i\,i+1}


M_i+


\frac{u_mu_n-u_1u_{m+1}}{u_{1\,n}}M_m


+\frac{u_1}{u_{1\,n}}\sum^{n-1}_{i=m+1}u_{i\,i+1}M_i,\quad


m=2,\dotsc,n-2,\\


&\dots\\


N_{n-1}&=N_n+\frac{u_n}{u_{1\,n}}\sum^{n-2}_{i=1}u_{i\,i+1}M_i+


\frac{u_nu_{n-1\,1}}{u_{1\,n}}M_{n-1}.


\end{split}


\end{equation}


Тогда система (11) с неизвестными $M_1,\dotsc,M_{n-1}$,


$N_n$ сводится к


\begin{equation*}


[\rho_1u_{2n}+(k^2-1)k_{1n}u_1u_{1\,2}]M_1+


\sum^{n-1}_{i=2}u_{i\,i+1}


[(k^2-1)k_{1n}u_1-\rho_1]M_i+


(k^2-1)k_{1n}u_{1n}N_n=0, \tag{17.1}


\end{equation*}


\begin{multline}


[(k^2-1)k_{2\,1}u_nu_{1\,2}-\rho_2u_{1n}]M_1+


[\rho_2u_{1n}+(k^2-1)k_{2\,1}u_1u_{2\,3}]M_2+\\


+(k^2-1)k_{2\,1}+u_1\sum^{n-1}_{i=3}u_{i\,i+1}M_i


+(k^2-1)k_{2\,1}u_{1n}N_n=0,


\tag{17.2}


\end{multline}


$$


\dots


$$


\begin{multline}


(k^2-1)k_{m\,m-1}u_n\sum^{m-2}_{i=1}u_{i\,i+1}M_i+


[(k^2-1)k_{m\,m-1}u_nu_{m-1\,m}+\rho_mu_{n1}]M_{m-1}+\\


+[\rho_mu_{1n}+(k^2-1)k_{m\,m-1}u_1u_{m\,m+1}]M_m+\\


+(k^2-1)k_{m\,m-1}u_1\sum^{n-1}_{i=m+1}u_{i\,i+1}M_i+


(k^2-1)k_{m\,m-1}u_{1n}N_n=0,\quad m=3,\dotsc,n-2,


\tag{17.$m$}


\end{multline}


$$


\dots


$$


\begin{multline*}


[\rho_n+(k^2-1)k_{n\,n-1}u_n]\sum^{n-2}_{i=1}u_{i\,i+1}M_i+


[\rho_nu_{n-1\,1}+(k^2-1)k_{n\,n-1}u_nu_{n-1\,n}]M_{n-1}+\\


+(k^2-1)k_{n\,n-1}u_{1n}N_n=0.


\end{multline*}


\addtocounter{equation}{-1}





Обозначим через $A$ матрицу системы (17). Для того чтобы эта система


однородных линейных уравнений имела нетривиальные решения, необходимо и


достаточно, чтобы


\begin{equation}


\det A=0.


\end{equation}


Приведя $A$ к треугольной форме, получим


\begin{equation}


A\sim


\begin{bmatrix}


N_n & M_{n-1} & \dots & M_1\\


b_{11} & b_{12} & \dots & b_{1n}\\


0 & b_{22} & \dots & b_{2n} \\


\hdotsfor[1.5]{4}\\


0 & 0 & & b_{nn}


\end{bmatrix}


.


\end{equation}


Отсюда следует, что система (17) имеет нетривиальные решения


тогда и только тогда, когда


\begin{equation}


b_{nn}=0.


\end{equation}





Таким образом, доказана





\begin{thmnonum}


Для того чтобы торообразная поверхность вращения с многоугольным


меридианом, имеющим вершины $A_i(u_i,\rho_i)$, $i=1,\dotsc,n$, была


нежесткой, необходимо


и достаточно, чтобы выполнялись условия {\em (18)} или эквивалентные им


условия {\em (20)}, где $A$ --- матрица системы {\em (17)} и


$u_{ij}=u_i-u_j$, $k_{ij}=k_i-k_j$,


$k_i=\dfrac{\rho_{i+1}-\rho_i}{u_{i+1}-u_i}$,


$k\ge 2$ --- целое число.


\end{thmnonum}





К.М.\,Белов \cite{Bel} исследовал торообразную поверхность вращения,


обладающую четырехугольным меридианом с вершинами


$A_1(-1,b)$, $A_2(0,b+c_1)$, $A_3(1,b)$, $A_4(0,b-c_2)$


способом, отличным от предлагаемого в этой статье, и получил соотношение


$$


\frac{1}{c_2}-\frac{1}{c_1}=\frac{k^2}{b},\quad


k\in\{2,3,\dotsc\},


$$


как необходимое и достаточное условие, чтобы такая поверхность была


нежесткой. То же соотношение можно получить из (18) или (20), используя


доказанную теорему.





Метод, использованный в данной статье, дает возможность определения поля


б.\,м. изгибаний. Из (19) при условии (20) получается редуцированная


система


\begin{align*}


b_{11}N_n+b_{12}M_{n-1}+\dotsb+b_{1n}M_1&=0,\\


b_{22}M_{n-1}+\dotsb+b_{2n}M_1&=0,\\


&\dots\\


b_{n-1\,n-1}M_2+b_{n-1\,n}M_1&=0,


\end{align*}


откуда (при условии $b_{n-1\,n-1}\ne 0$) находим


$M_2=-\dfrac{b_{n-1\,n}}{b_{n-1\,n-1}}M_1$ и далее выражаем


$M_3,\dotsc,M_{n-1}$,


$N_n$ через $M_1$


(неопределенная постоянная). Из (14) и (16), находим выражение


$M_n$,


$N_1,N_2,\dotsc,N_{n-1}$, через $M_1$. Подставляя найденные выражения в


(6), определяем $\psi_i(u)$ и, тем самым, поле б.\,м. изгибаний.





\begin{thebibliography}{9}





\bibitem{Bel}


Белов К.М. {\em О бесконечно малых изгибаниях торообразующей поверхности


вращения} // Сиб. матем. журн. -- 1968. -- Т.\,9. -- \No\,3. --


С.\,490--494.


\bibitem{Kohn}


Кон-Фоссен С.Э. {\em Некоторые вопросы дифференциальной геометрии в


целом}. -- М.: Физматгиз, 1959.








\end{thebibliography}





\vskip 2cm





\post{\it Университет г.\,Нише}{\it Поступила}


\post{({\it Югославия})}{11.05.1995}





\end{document}


